% =============================================================================
% Kapitel 5 -- Zusammenfassung und Fazit
% =============================================================================
\section{Zusammenfassung und Fazit}

\subsection{Ergebnisse}

CardSiege ist ein vollständiges, spielbares Kartenspiel, das alle funktionalen
Anforderungen erfüllt und neun GoF-Design-Patterns sinnvoll integriert. Die vier
Erzeugungsmuster -- Singleton, Factory Method, Prototype und Builder -- strukturieren
gemeinsam die gesamte Objekterzeugung des Spiels: Der Singleton sichert den
konsistenten Spielzustand, die Factory entkoppelt das Deck von konkreten
Kartentypen, das Prototype-Pattern ermöglicht effizientes Klonen konfigurierter
Objekte, und der Builder trennt Turmkonfiguration von Platzierung. Das
Decorator-Pattern als einziges Strukturmuster löst das zentrale Erweiterungsproblem
im Turmsystem. Die vier Verhaltensmuster -- Command, Strategy, State und Observer --
modellieren das laufende Spielgeschehen so, dass Erweiterungen ohne Modifikation
bestehender Klassen möglich sind.

\subsection{Erkenntnisse und Grenzen}

Die Arbeit hat gezeigt, dass Design Patterns erst im Kontext eines realen Problems
ihren vollen Wert entfalten. Das Decorator-Pattern wirkt in der Theorie
überengineert -- erst das Skizzieren der alternativen Vererbungshierarchie macht
den konkreten Mehrwert sichtbar. Das Strategy-Pattern reduziert eine potentiell
komplexe Fallunterscheidung im Turmangriff auf zwei austauschbare Algorithmen,
die sich ohne Seiteneffekte kombinieren lassen.

Besonders aufschlussreich war das natürliche Zusammenspiel der Patterns: Factory
und Prototype kooperieren im Deck, State und Observer im Phasenzyklus, Command
und Decorator im Buff-Kartensystem. Die Patterns verstärken sich gegenseitig --
keine Lösung ist in sich geschlossen, alle greifen ineinander.

Grenzen wurden ebenfalls sichtbar. Das Singleton erschwert Unit-Tests durch
persistierenden Zustand. Der \texttt{GameManager} trägt trotz aller Delegation
noch sehr viele Verantwortlichkeiten; in einem größeren Projekt würde man ihn in
separate \texttt{CombatManager}- und \texttt{ResourceManager}-Klassen aufteilen.

\subsection{Ausblick}

Das System bietet konkrete Ansatzpunkte für Erweiterungen, die unmittelbar auf
der bestehenden Architektur aufbauen. Konfigurierbare Decks und ein
Progressionssystem, bei dem nach und nach neue Karten freigeschaltet werden,
ließen sich vollständig über das Factory-Pattern realisieren: Jede
Fortschrittsstufe entspräche einer neuen \texttt{CardFactory}-Implementierung,
ohne eine Zeile in \texttt{Deck} oder den bestehenden Karten zu ändern. Neue
Gegner- und Turmtypen sind durch \texttt{EnemyPrototype} beziehungsweise
\texttt{StrategyCardFactory} bereits architektonisch vorbereitet -- das
Hinzufügen eines weiteren Eintrags genügt.

Das Command-Pattern legt die Grundlage für eine Undo-Funktion: Eine
\texttt{undo()}-Methode auf \texttt{Card} würde Kartenaktionen rückgängig
machbar machen, was gerade für ein Progressionssystem mit begrenzten Ressourcen
spielerisch wertvoll wäre. Das Memento-Pattern böte eine elegante Basis für
Spielstand-Speicherung. Ein Endless-Mode schließlich ließe sich mit minimalen
Eingriffen umsetzen: \texttt{WaveSpawner} und \texttt{GameOverPhase} müssten
lediglich die feste Wellenobergrenze durch eine unbegrenzte Skalierungslogik
ersetzen, alle anderen Systeme blieben unverändert.
